\documentclass{article}
\usepackage{amsmath,amssymb,amsthm}
\newtheorem{assumption}{Assumption}
\newtheorem{definition}{Definition}
\newtheorem{theorem}{Theorem}
\newtheorem{lemma}{Lemma}
\newtheorem{proposition}{Proposition}
\newtheorem{openendedquestion}{Open-ended Question}
\newtheorem{conjecture}{Conjecture}
\newcommand{\coreref}[1]{\ref{#1}}

\begin{document}

\section*{stat_prte_ot_boundary_adaptation (v1)}

\paragraph{TL;DR.} The proposal studies the sharp lower endpoint of the continuous-instrument CCOT PRTE identified interval when inference is driven by observations whose realized propensity sits close to the lower support boundary. The formal kernel is a boundary-local decomposition of the empirical CCOT plug-in endpoint into a regular interior influence term, a truncated shell score, and an explicit extreme-shell residual from 0 < B <= n^(-1), all weighted by chi_P(O) B^(-1/2), with the tail exponent alpha determining the normalization rho_n(alpha) = sqrt(n / K_alpha(n^(-1))). The boundary class pins down the shell law with exact regular variation P(0 < B <= t) = c_B t^alpha {1 + o(1)}, atom-free boundary support, stabilization of the boundary marks (chi_P, active-face label), centered shell nondegeneracy, and explicit Gaussian or Levy limit objects for the truncated and extreme shell coordinates. The phase transition sits at alpha = 1: alpha > 1 yields a root-n Gaussian endpoint law, alpha = 1 yields a distinct logarithmic transition law, and 0 < alpha < 1 yields a non-Gaussian boundary-irregular limit in which the extreme shell enters at first order. The stochastic conclusions are now stated honestly as conditional on a high-level active-face local linearization of the CCOT endpoint together with theorem-level empirical-path admissibility at the boundary normalization scale, rather than as a derivation of that singular derivative from the CCOT primitives. Relative to the in-repository note stat_ate_overlap_decay, which studies overlap-tail phase behavior for ATE risk, this artifact is about a sharp CCOT PRTE endpoint law with an explicit shell-process decomposition and an endpoint-level phase transition. TARGET ADJUSTED (D0): the field contribution is the positive side only — the alpha-indexed boundary-local endpoint limit law with the alpha*=1 phase transition (thm:endpoint-limit-law, thm:local-experiment, prop:phase-normalization, prop:boundary-nonuniformity), the regular root-n strict-overlap reduction (prop:strict-overlap-reduction), and the fixed-alpha honest CI as a conditional open-ended result (oeq:boundary-honest-ci, conditional on a fixed-alpha calibrator assumption that is not constructible from primitives). No cross-alpha honest-adaptation theorem is claimed in this core.

\section{Project justification}

\paragraph{Gap.} TanBlanchetSyrgkanis2026PRTEOT gives the sharp CCOT PRTE lower endpoint and a continuous-instrument convergence rate, but not the endpoint normalization, marked boundary-shell experiment, or CI theory. SasakiUra2018PRTEInference covers regular point-identified PRTE inference, while ImbensManski2004PartialCI and Stoye2009MorePartialCI cover regular interval endpoints. Within the repository, stat_ate_overlap_decay develops overlap-tail phase bookkeeping for ATE risk rather than a CCOT PRTE endpoint law or a fixed-alpha endpoint CI problem.

\paragraph{Niche.} The niche is the lower-propensity boundary regime where the endpoint derivative is carried by a unique active transport face and the relevant shell law can be summarized by one exponent alpha, a limiting boundary-mark distribution M_P, and phase-specific covariance or Levy objects. In that regime the inferential object is neither the regular PRTE score of Sasaki-Ura nor a generic moment-inequality bound endpoint; it is a boundary-local OT endpoint with a phase transition at alpha = 1.

\paragraph{Fill.} The project fills that niche with a theorem-level package: an explicit alpha-indexed normalization, a truncated-domain local experiment with identified shell covariance or joint Levy structure, an endpoint limit law that keeps the extreme shell explicit while separating regular and boundary-irregular phases, a strict-overlap reduction to the regular partial-ID CI template, and a conditional fixed-alpha CI handle that records the remaining calibration problem without overstating a cross-alpha adaptation theorem. The CCOT-specific singular derivative itself is treated honestly as a provisional high-level active-face expansion input, and the empirical plug-in path is controlled only through a separate theorem-level admissibility/rate assumption rather than through a claimed primitive Donsker derivation inside this core.

\section{Related work}

TanBlanchetSyrgkanis2026PRTEOT is the anchor paper because it expresses sharp PRTE bounds as a constrained conditional optimal-transport problem and proves a continuous-instrument convergence rate for the lower endpoint. The present kernel keeps that exact endpoint but adds the missing normalization, marked boundary-shell experiment, and CI-facing limit law with the omitted extreme shell made explicit. This core does not claim to derive the singular active-face endpoint linearization from the CCOT primitives themselves; it treats that step as a high-level interface assumption and then studies the stochastic consequences. The closest in-repository comparator is stat_ate_overlap_decay: it studies one-sided overlap-tail phase behavior for ATE risk and minimax bookkeeping, whereas the present artifact keeps the CCOT PRTE endpoint fixed and derives a law-level phase transition for that endpoint itself. SasakiUra2018PRTEInference is the closest PRTE inference comparator because it derives regular asymptotics for point-identified PRTEs with generated components; it is informative for the interior influence term phi_P, but it does not treat a sharp OT endpoint with boundary-local irregularity. DaouiaFlorensSimar2010FrontierEVT is the nearest boundary-law analogue because it shows how support-edge localization can generate nonstandard limits and stable-law behavior, but its frontier object is not an active-face OT value. FangSantos2019DirectionalInference is the main nonsmooth-inference template because it explains why ordinary bootstrap logic can fail once the endpoint is only directionally regular; here it is guidance for the downstream limit-law consequences rather than a claimed source of the CCOT linearization itself. FranguridiLiu2025RegularizedOTPartialID is the closest OT inference paper because it smooths a partially identified OT value and resamples the regularized functional, whereas the present proposal keeps the unregularized CCOT endpoint and studies the truncated-shell and extreme-shell singularities directly. ImbensManski2004PartialCI and Stoye2009MorePartialCI are the exact regular partial-ID CI comparators; prop:strict-overlap-reduction is the formal bridge back to that benchmark when the lower propensity stays away from the external boundary. CaiLow2005AdaptiveCI and RobinsVanderVaart2006AdaptiveSets remain background adaptation references only: they explain why cross-class honest-length questions are natural, but no PRTE-specific honest-adaptation theorem is claimed in this core. KimPollard1990CubeRoot remains an adjacent proof-discipline guide because it is the canonical reminder that localized nonsmooth criteria can have non-Gaussian limits even when global Gaussian heuristics look tempting.

\section{Setup}

Target estimand: theta_L = inf_{pi in Pi_ccot(P)} Int m_P(u,v,x) d pi(u,v,x).
Estimand-defining functional: hat_theta_L_n = Gamma(hat_P_n), with rho_n(alpha){hat_theta_L_n - theta_L} = rho_n(alpha) n^(-1) sum_i phi_P(O_i) + S_nP + R_nP + o_P(1) and rho_n(alpha) = sqrt(n / K_alpha(n^(-1))).
\begin{itemize}
  \item $n$ --- sample_size; N; index $:= Sample size.$
  \item $O$ --- random_vector; {(Y,D,Z,X): Y in R, D in {0,1}, Z in Z, X in X}; data_unit $:= Observed unit O = (Y, D, Z, X).$
  \item $O_1, ..., O_n$ --- sample; ({(Y,D,Z,X): Y in R, D in {0,1}, Z in Z, X in X})^n; data $:= Observed sample of n units.$
  \item $Y$ --- outcome; R; data_component $:= Observed outcome component of O.$
  \item $D$ --- treatment; {0,1}; data_component $:= Observed treatment indicator component of O.$
  \item $Z$ --- instrument; Z; data_component $:= Observed continuous instrument component of O.$
  \item $X$ --- covariate; X; data_component $:= Observed covariate component of O.$
  \item $P$ --- law; laws of O; model $:= Observed-data law of O.$
  \item $p$ --- propensity_map; primitive_function $:= p(z,x) = P(D = 1 | Z = z, X = x).$
  \item $underline_p$ --- support_boundary_map; boundary_object $:= underline_p(x) = inf_{z in supp(Z | X = x)} p(z,x).$
  \item $B$ --- boundary_radius; [0,1]; boundary_localization $:= B = p(Z,X) - underline_p(X), the lower-propensity boundary distance of the realized instrument value.$
  \item $Pi_ccot(P)$ --- transport_feasible_set; set of conditional couplings indexed by X; identified_set_primitive $:= Feasible conditional optimal-transport couplings entering the sharp PRTE lower bound under P.$
  \item $m_P$ --- transport_payoff_kernel; identified_set_primitive $:= PRTE payoff kernel in the conditional optimal-transport representation under P.$
  \item $Gamma$ --- endpoint_map; endpoint_map $:= Gamma(P) is the lower endpoint map of the sharp CCOT PRTE identified interval.$
  \item $theta_L$ --- partial_identification_endpoint; R; target_estimand $:= theta_L = Gamma(P), the lower endpoint of the sharp PRTE identified interval.$
  \item $A$ --- active_face_label_space; finite or countable label set; geometry_index $:= Index set for active CCOT transport faces near the lower propensity boundary.$
  \item $a_P$ --- active_face_label_map; boundary_geometry $:= a_P(O) is the label of the active lower-endpoint CCOT transport face at observation O on the region 0 < B <= bar_b.$
  \item $chi_P$ --- boundary_loading; boundary_derivative $:= Boundary loading in the directional derivative of Gamma at P along the unique active transport face.$
  \item $phi_P$ --- influence_function; regular_component $:= Regular interior influence function for perturbations of Gamma that stay away from the boundary-local region B = 0.$
  \item $bar_b$ --- boundary_cutoff; (0,1); localization_constant $:= Fixed upper cutoff delimiting the boundary-local region 0 < B <= bar_b.$
  \item $alpha$ --- tail_exponent; (0, infinity); regime_parameter $:= Boundary-tail exponent governing the mass of laws with small boundary distance B.$
  \item $c_B$ --- tail_scale_constant; (0, infinity); tail_constant $:= Exact regular-variation constant in P(0 < B <= t) = c_B t^alpha {1 + o(1)}.$
  \item $M_P$ --- mark_limit_law; probability measures on R x A; shell_mark_distribution $:= Weak limit of the conditional law of (chi_P(O), a_P(O)) given 0 < B <= t as t downarrow 0.$
  \item $K_alpha$ --- tail_scale_function; normalization_object $:= Boundary variance scale K_alpha(t) = 1 if alpha > 1, K_alpha(t) = 1 + log(1 / t) if alpha = 1, and K_alpha(t) = t^(alpha - 1) if 0 < alpha < 1.$
  \item $rho_n(alpha)$ --- normalization_sequence; (0, infinity)^N; rate_normalization $:= rho_n(alpha) = sqrt(n / K_alpha(n^(-1))).$
  \item $hat_P_n$ --- empirical_law; random laws on O; empirical_object $:= hat_P_n = n^(-1) sum_{i=1}^n delta_{O_i}.$
  \item $hat_theta_L_n$ --- plugin_endpoint_estimator; R; estimator $:= hat_theta_L_n = Gamma(hat_P_n), the empirical CCOT plug-in lower endpoint.$
  \item $G_n,P,u$ --- shell_summand_map; local_summand $:= G_n,P,u(O) = rho_n(alpha) chi_P(O) B^(-1/2) 1{ u n^(-1) < B <= bar_b }.$
  \item $J_n,P$ --- extreme_shell_summand_map; extreme_shell_summand $:= J_n,P(O) = rho_n(alpha) chi_P(O) B^(-1/2) 1{ 0 < B <= n^(-1) }.$
  \item $S_nP$ --- boundary_score; R; boundary_statistic $:= S_nP = n^(-1) sum_{i=1}^n [ G_n,P,1(O_i) - E_P G_n,P,1(O) ].$
  \item $R_nP$ --- extreme_shell_statistic; R; extreme_shell_residual $:= R_nP = n^(-1) sum_{i=1}^n [ J_n,P(O_i) - E_P J_n,P(O) ].$
  \item $u_0$ --- truncation_index; (0,1]; process_domain_parameter $:= Fixed lower shell index used to avoid the singular point u = 0.$
  \item $H_nP_u0$ --- truncated_boundary_process; local_experiment_process $:= For u in [u_0,1], H_nP_u0(u) = n^(-1) sum_{i=1}^n [ G_n,P,u(O_i) - E_P G_n,P,u(O) ].$
  \item $Sigma_P,u0$ --- covariance_kernel; limit_covariance $:= Positive-semidefinite covariance kernel governing the Gaussian or logarithmic finite-dimensional shell limits on [u_0,1].$
  \item $Lambda_P,alpha,u0$ --- levy_measure_family; family of Levy measures on finite-dimensional vectors with one extreme-shell coordinate and truncated-shell coordinates indexed by [u_0,1]; jump_limit_object $:= Levy-measure family governing the joint infinitely divisible limits of the extreme-shell coordinate J_n,P and the truncated shell coordinates G_n,P,u when 0 < alpha < 1.$
  \item $H_P,alpha,u0$ --- limit_process; D([u_0,1]); limit_object $:= Boundary-local limit experiment for H_nP_u0 on [u_0,1]: when alpha >= 1 it is the Gaussian process with covariance kernel Sigma_P,u0, and when 0 < alpha < 1 it is the unique continuous-path law on D([u_0,1]) whose finite-dimensional marginals are the truncated-shell coordinates of the infinitely divisible vectors indexed by Lambda_P,alpha,u0.$
  \item $p_so$ --- strict_overlap_floor; (0,1); sanity_parameter $:= Deterministic external lower-propensity floor for the strict-overlap comparator class.$
  \item $alpha_0, alpha_1$ --- tail_exponents; (0, infinity)^2; adaptation_parameters $:= Two tail exponents used in the between-class honest-CI adaptation comparison.$
  \item $delta$ --- coverage_level; (0,1); inference_parameter $:= Target miscoverage level for confidence intervals.$
  \item $C_n$ --- confidence_interval_sequence; lower_bound_object $:= Generic confidence interval sequence placeholder for hypothetical cross-class honesty questions about theta_L.$
  \item $C_n_alpha$ --- confidence_interval_sequence; open_problem_object $:= Prospective alpha-specific confidence interval sequence for theta_L posed in oeq:boundary-honest-ci.$
\end{itemize}

\section{Assumptions}

\begin{assumption}[ass:iid]\label{ass:iid}
O_1, ..., O_n are i.i.d. draws from P.
\par\smallskip\noindent\emph{Standard} (i.i.d. sampling, \cite{TanBlanchetSyrgkanis2026PRTEOT}).
\end{assumption}

\begin{assumption}[ass:boundary-atom-free]\label{ass:boundary-atom-free}
P(B = 0) = 0.
\par\smallskip\noindent\emph{Novel.} This is realistic when the generalized propensity map reaches its lower support boundary only along a null set of instrument realizations, so the continuous-IV design has a thin boundary neighborhood rather than a positive-mass pile-up at the edge.
\end{assumption}

\begin{assumption}[ass:boundary-tail]\label{ass:boundary-tail}
There exists c_B > 0 such that P(0 < B <= t) = c_B t^alpha {1 + o(1)} as t downarrow 0.
\par\smallskip\noindent\emph{Novel.} This holds in latent-index IV designs where the realized propensity distance from the lower face has an exact regularly varying tail. It is the natural primitive version of the alpha-indexed boundary-mass regime singled out by the CCOT PRTE endpoint.
\end{assumption}

\begin{assumption}[ass:boundary-mark-stabilization]\label{ass:boundary-mark-stabilization}
There exists a probability measure M_P on R x A such that the conditional law of (chi_P(O), a_P(O)) given 0 < B <= t converges weakly to M_P as t downarrow 0.
\par\smallskip\noindent\emph{Novel.} This is realistic when the active face and its derivative loading stabilize along the lower propensity edge, which is exactly the regime where a boundary-local asymptotic experiment should exist. It permits the shell law to be summarized by a limiting mark distribution instead of by ad hoc shell-wise constants.
\end{assumption}

\begin{assumption}[ass:centered-shell-variance]\label{ass:centered-shell-variance}
There exist c_V_minus > 0 and c_V_plus > 0 such that c_V_minus K_alpha(t) <= Var_P( chi_P(O) B^(-1/2) 1{ t < B <= bar_b } ) <= c_V_plus K_alpha(t) for every t in (0, bar_b].
\par\smallskip\noindent\emph{Novel.} This holds when the boundary loading is not asymptotically canceled by sign changes within shrinking shells, so the singular shell contributes at the alpha-indexed order identified by K_alpha. It is the centered nondegeneracy condition needed for a genuine phase transition rather than a variance upper bound alone.
\end{assumption}

\begin{assumption}[ass:unique-active-face]\label{ass:unique-active-face}
For P-almost every X and every realized boundary point with 0 < B <= bar_b, the lower-endpoint CCOT program has a unique active transport face.
\par\smallskip\noindent\emph{Novel.} This is realistic on generic subclasses where the supporting hyperplane to the CCOT value is unique except on tie sets of codimension one. It isolates the identifiable geometry already implicit in the anchor paper's one-dimensional transport formulas and keeps the proposal off the tie/slack frontier.
\end{assumption}

\begin{assumption}[ass:dual-gap-scale]\label{ass:dual-gap-scale}
There exists c_gap > 0 such that, on the region 0 < B <= bar_b, the dual value gap between the active face and the closest competing face is at least c_gap B.
\par\smallskip\noindent\emph{Novel.} This is realistic in smooth support-price models where the winning transport face remains separated from runners-up at first order as one approaches the lower propensity edge. It is the exact genericity condition needed to keep the endpoint derivative from being dominated by accidental face switches.
\end{assumption}

\begin{assumption}[ass:gaussian-shell-lindeberg]\label{ass:gaussian-shell-lindeberg}
If alpha >= 1, then for every fixed u_0 in (0,1] there exists a continuous positive-semidefinite kernel Sigma_P,u0 on [u_0,1]^2 such that for every finite set u_1, ..., u_m in [u_0,1], n^(-1) Cov_P( G_n,P,u_j(O), G_n,P,u_k(O) ) -> Sigma_P,u0(u_j,u_k) for all j,k and the triangular array { n^(-1/2) [ G_n,P,u_j(O_i) - E_P G_n,P,u_j(O) ] : i <= n } satisfies the multivariate Lindeberg condition.
\par\smallskip\noindent\emph{Novel.} This is realistic as a high-level shell CLT premise on subclasses where the boundary-score array is generated by many small shell contributions with uniformly negligible single-observation impact after truncation away from u = 0. It is imposed only to encode the Gaussian finite-dimensional limit object on fixed truncated domains, not as a claimed published condition for the CCOT PRTE endpoint.
\end{assumption}

\begin{assumption}[ass:stable-shell-domain]\label{ass:stable-shell-domain}
If 0 < alpha < 1, then for every fixed u_0 in (0,1] and every finite set u_1, ..., u_m in [u_0,1], there exists a Levy measure Lambda_P,alpha,u0 on R^(m+1) \ {0} such that n P( ( n^(-1) J_n,P(O), n^(-1) G_n,P,u_1(O), ..., n^(-1) G_n,P,u_m(O) ) in E ) -> Lambda_P,alpha,u0(E) for every Borel set E bounded away from zero, and the centered array has negligible Gaussian part.
\par\smallskip\noindent\emph{Novel.} This is realistic as a high-level irregular-shell premise when the endpoint is driven by rare boundary observations whose rescaled contributions converge to a Poissonian or stable jump law on every fixed truncated shell domain. It is imposed only to encode the non-Gaussian finite-dimensional limit object for the CCOT PRTE shell array, not as a textbook condition taken verbatim from frontier extreme-value theory.
\end{assumption}

\begin{assumption}[ass:regular-linearization]\label{ass:regular-linearization}
For perturbations of P supported on the interior region { B > bar_b / 2 }, the endpoint map Gamma is pathwise differentiable with influence function phi_P in L2(P).
\par\smallskip\noindent\emph{Novel.} This is realistic because once the perturbation is separated from the propensity boundary the CCOT endpoint behaves like a regular smooth functional of the observed law, mirroring the orthogonal PRTE asymptotics in Sasaki-Ura. The proposal uses it only to isolate the regular interior term from the boundary-local singular component.
\end{assumption}

\begin{assumption}[ass:strict-overlap-propensity]\label{ass:strict-overlap-propensity}
underline_p(X) >= p_so almost surely.
\par\smallskip\noindent\emph{Novel.} This is realistic in continuous-IV designs where the generalized propensity image is externally trimmed away from zero because institutional or design restrictions prevent instruments from pushing any covariate stratum arbitrarily close to the treatment boundary. It is used only for the regular strict-overlap comparator subclass.
\end{assumption}

\begin{assumption}[ass:strict-overlap-boundary-regularity]\label{ass:strict-overlap-boundary-regularity}
There exists C_so < infinity such that E_P[ chi_P(O)^2 B^(-1) 1{ 0 < B <= t } ] <= C_so for every t in (0, bar_b].
\par\smallskip\noindent\emph{Novel.} This is realistic in strict-overlap submodels where the external propensity floor removes the thin-tail mechanism and the active-face loading remains square-integrable all the way to the internal support edge. It is exactly the regular-endpoint comparator needed for the sanity reduction and is not imposed on the main alpha-tail class.
\end{assumption}

\begin{assumption}[ass:shell-row-lindeberg]\label{ass:shell-row-lindeberg}
For alpha>=1, every fixed u_0 in (0,1] and every finite u_1,...,u_m in [u_0,1], the row variables n^(-1)(G_n,P,u_j(O_i)-E_P G_n,P,u_j(O))_{j=1}^m satisfy the multivariate Lindeberg condition, and their row-sum covariance n^(-1)Cov_P(G_n,P,u_j,G_n,P,u_k) converges to Sigma_P,u0(u_j,u_k).
\par\smallskip\noindent\emph{Novel.} standard: Lindeberg-Feller triangular-array condition at the row-summand scale
\end{assumption}

\begin{assumption}[ass:shell-increment-moment]\label{ass:shell-increment-moment}
For every fixed u_0 in (0,1], there exist q>2, beta>1, and C<infinity such that sup_n E_P|H_nP_u0(v)-H_nP_u0(u)|^q <= C|v-u|^beta for all u,v in [u_0,1], and sup_n E_P|H_nP_u0(u_0)|^q <= C.
\par\smallskip\noindent\emph{Novel.} standard: Kolmogorov-Chentsov/Billingsley-type moment tightness criterion
\end{assumption}

\begin{assumption}[ass:stable-shell-drift]\label{ass:stable-shell-drift}
If 0<alpha<1, then for every fixed u_0 and finite u_1,...,u_m in [u_0,1], the compensated drift terms for the centered row array n^(-1)(J_n,P,G_n,P,u_1,...,G_n,P,u_m) converge for a common bounded truncation function.
\par\smallskip\noindent\emph{Novel.} standard: drift condition in the triangular-array Levy convergence theorem
\end{assumption}

\begin{assumption}[ass:stable-shell-nondegenerate-jump]\label{ass:stable-shell-nondegenerate-jump}
If 0<alpha<1, the projection of Lambda_P,alpha,u0 onto the extreme-shell coordinate is nonzero for every fixed u_0 in (0,1].
\par\smallskip\noindent\emph{Novel.} novel: a no-cancellation/nondegeneracy condition for the CCOT boundary loading
\end{assumption}

\begin{assumption}[ass:extreme-shell-first-moment]\label{ass:extreme-shell-first-moment}
For alpha>=1, there are C<infinity and t_0>0 such that E_P[|chi_P(O)| B^(-1/2) 1{0<B<=t}] <= C t^(alpha-1/2) for all 0<t<=t_0.
\par\smallskip\noindent\emph{Novel.} standard: local first-moment envelope, implied for example by bounded boundary loading and a regularly varying boundary density
\end{assumption}

\begin{assumption}[ass:boundary-annulus-oscillation]\label{ass:boundary-annulus-oscillation}
When alpha=1, for every sufficiently small delta_0>0 there are sequences u_n,v_n in (0,delta_0] and constants eta,epsilon>0 such that liminf_n P_P(|H_nP_delta0(u_n)-H_nP_delta0(v_n)|>eta)>=epsilon.
\par\smallskip\noindent\emph{Novel.} novel: boundary-annulus no-cancellation condition at the logarithmic transition
\end{assumption}

\begin{assumption}[ass:regular-boundary-joint-limit]\label{ass:regular-boundary-joint-limit}
If alpha>1, the joint row array formed by n^(-1/2)phi_P(O_i) and n^(-1)(G_n,P,u_j(O_i)-E_PG_n,P,u_j(O)) for every finite u_1,...,u_m satisfies Lindeberg and has a finite limiting covariance matrix.
\par\smallskip\noindent\emph{Novel.} standard: joint Lindeberg condition for regular and boundary empirical arrays
\end{assumption}

\begin{assumption}[ass:fixed-alpha-calibrator]\label{ass:fixed-alpha-calibrator}
For each fixed alpha, the data-driven law estimator hat L_n,alpha in the corrected CI construction satisfies sup_{P in P_n(alpha)} sup_t |hat L_n,alpha((-infinity,t])-P_P(T_n,alpha<=t)| -> 0 in probability at continuity points of the limiting law, its delta/2 and 1-delta/2 quantiles are uniformly consistent, and E_P[hat q_+(alpha)-hat q_-(alpha)] is uniformly O(1).
\par\smallskip\noindent\emph{Novel.} standard: high-level bootstrap/plug-in calibration consistency condition
\end{assumption}

\begin{assumption}[ass:strict-overlap-empirical-smoothness]\label{ass:strict-overlap-empirical-smoothness}
For each P in P_n^so(p_so) there exist a class F_P of square-integrable measurable functions, a seminorm ||Q-P||_{F_P}=sup_{f in F_P}|(Q-P)f| on signed perturbations of P, and a deterministic modulus eta_P(s) with eta_P(s) -> 0 as s downarrow 0 such that phi_P belongs to F_P, E_P phi_P(O)=0, E_P phi_P(O)^2<infinity, ||hat_P_n-P||_{F_P}=O_P(n^(-1/2)), and for every law Q in a strict-overlap neighborhood of P, |Gamma(Q)-Gamma(P)-(Q-P)phi_P| <= ||Q-P||_{F_P} eta_P(||Q-P||_{F_P}).
\par\smallskip\noindent\emph{Novel.} standard: local Hadamard differentiability plus a Donsker or stochastic-equicontinuity empirical-process bound for a regular plug-in functional
\end{assumption}

\begin{assumption}[ass:regular-influence-centering]\label{ass:regular-influence-centering}
The regular influence function is chosen in its canonical centered version: E_P phi_P(O)=0 and E_P phi_P(O)^2<infinity; if a pathwise derivative representative is uncentered, replace it by phi_P-P phi_P without changing (hat P_n-P)phi_P.
\par\smallskip\noindent\emph{Novel.} standard: canonical normalization of an influence function in semiparametric theory
\end{assumption}

\begin{assumption}[ass:irregular-endpoint-projection-nondegeneracy]\label{ass:irregular-endpoint-projection-nondegeneracy}
If 0<alpha<1, then for every fixed u_0 in (0,1], the projection of Lambda_P,alpha,u0 under the scalar endpoint map that sends the extreme-shell and u=1 truncated-shell coordinates to R+H(1) is a nonzero Levy measure after deterministic centering.
\par\smallskip\noindent\emph{Novel.} novel: no-cancellation condition for the CCOT boundary endpoint projection
\end{assumption}

\begin{assumption}[ass:ccot-gamma-hadamard]\label{ass:ccot-gamma-hadamard}
There exist a local active-face neighborhood U_P of P, a seminorm ||.||_F on signed law perturbations, the measurable boundary score f_B(O)=chi_P(O)B(O)^(-1/2)1{0<B(O)<=bar_b}, and a deterministic modulus eta_P(s) with eta_P(s) -> 0 as s downarrow 0 such that for every law Q in U_P, Gamma(Q)-Gamma(P)=(Q-P)phi_P+(Q-P)f_B+Rem_P(Q) and |Rem_P(Q)| <= ||Q-P||_F eta_P(||Q-P||_F).
\par\smallskip\noindent\emph{Novel.} D0.R provisional rescope: this is the high-level active-face local linearization input actually used by the endpoint-law proof. The core no longer presents the singular boundary derivative as a standard CCOT primitive fact derived inside the note; instead it conditions on that expansion and studies the stochastic consequences once it is available.
\end{assumption}

\begin{assumption}[ass:ccot-donsker]\label{ass:ccot-donsker}
With the neighborhood U_P, seminorm ||.||_F, and modulus eta_P from ass:ccot-gamma-hadamard, P(hat_P_n in U_P) -> 1 and ||hat_P_n-P||_F = O_P(rho_n(alpha)^(-1)).
\par\smallskip\noindent\emph{Novel.} D0.R provisional rescope: despite the legacy id, this node is not a discharged Donsker CLT claim. The endpoint proof only uses neighborhood entry and boundary-normalization control of the empirical path, so this assumption records exactly the theorem-needed empirical admissibility/rate input required to turn the local active-face expansion into an o_P(rho_n(alpha)^(-1)) remainder.
\end{assumption}

\section{Definitions}

\begin{definition}[def:ccot-lower-endpoint]\label{def:ccot-lower-endpoint}
$Gamma(P) := Gamma(P) = inf_{pi in Pi_ccot(P)} Int m_P(u,v,x) d pi(u,v,x)$ \quad (\text{inputs: } Pi_ccot(P), m_P, P)
\end{definition}

\begin{definition}[def:boundary-class]\label{def:boundary-class}
$P_n(alpha) := { P : P satisfies ass:boundary-atom-free, ass:boundary-tail, ass:boundary-mark-stabilization, ass:centered-shell-variance, ass:unique-active-face, ass:dual-gap-scale, ass:regular-linearization, ass:shell-row-lindeberg when alpha >= 1, and ass:stable-shell-domain when 0 < alpha < 1 }$ --- the class of all objects satisfying: ass:boundary-atom-free; ass:boundary-tail; ass:boundary-mark-stabilization; ass:centered-shell-variance; ass:unique-active-face; ass:dual-gap-scale; ass:shell-row-lindeberg; ass:stable-shell-domain; ass:regular-linearization.
\end{definition}

\begin{definition}[def:strict-overlap-class]\label{def:strict-overlap-class}
$P_n^so := { P : P satisfies ass:strict-overlap-propensity, ass:strict-overlap-boundary-regularity, ass:unique-active-face, ass:dual-gap-scale, and ass:regular-linearization }$ --- the class of all objects satisfying: ass:strict-overlap-propensity; ass:strict-overlap-boundary-regularity; ass:unique-active-face; ass:dual-gap-scale; ass:regular-linearization.
\end{definition}

\begin{definition}[def:tail-scale]\label{def:tail-scale}
$K_alpha := K_alpha(t) = 1 if alpha > 1, K_alpha(t) = 1 + log(1 / t) if alpha = 1, and K_alpha(t) = t^(alpha - 1) if 0 < alpha < 1$ \quad (\text{inputs: } alpha, t)
\end{definition}

\begin{definition}[def:alpha-normalization]\label{def:alpha-normalization}
$rho_n(alpha) := rho_n(alpha) = sqrt(n / K_alpha(n^(-1)))$ \quad (\text{inputs: } n, K_alpha, alpha)
\end{definition}

\begin{definition}[def:empirical-law]\label{def:empirical-law}
$hat_P_n := hat_P_n = n^(-1) sum_{i=1}^n delta_{O_i}$ \quad (\text{inputs: } O_1, ..., O_n, n)
\end{definition}

\begin{definition}[def:plugin-endpoint]\label{def:plugin-endpoint}
$hat_theta_L_n := hat_theta_L_n = Gamma(hat_P_n)$ \quad (\text{inputs: } Gamma, hat_P_n)
\end{definition}

\begin{definition}[def:shell-summand]\label{def:shell-summand}
$G_n,P,u := For u in (0,1], G_n,P,u(O) = rho_n(alpha) chi_P(O) B^(-1/2) 1{ u n^(-1) < B <= bar_b }$ \quad (\text{inputs: } rho_n(alpha), chi_P, B, bar_b, n, P)
\end{definition}

\begin{definition}[def:extreme-shell-summand]\label{def:extreme-shell-summand}
$J_n,P := J_n,P(O) = rho_n(alpha) chi_P(O) B^(-1/2) 1{ 0 < B <= n^(-1) }$ \quad (\text{inputs: } rho_n(alpha), chi_P, B, n, P)
\end{definition}

\begin{definition}[def:boundary-score]\label{def:boundary-score}
$S_nP := S_nP = n^(-1) sum_{i=1}^n [ G_n,P,1(O_i) - E_P G_n,P,1(O) ]$ \quad (\text{inputs: } n, G_n,P,u, P, O_1, ..., O_n)
\end{definition}

\begin{definition}[def:extreme-shell-residual]\label{def:extreme-shell-residual}
$R_nP := R_nP = n^(-1) sum_{i=1}^n [ J_n,P(O_i) - E_P J_n,P(O) ]$ \quad (\text{inputs: } n, J_n,P, P, O_1, ..., O_n)
\end{definition}

\begin{definition}[def:truncated-boundary-process]\label{def:truncated-boundary-process}
$H_nP_u0 := For u in [u_0,1], H_nP_u0(u) = n^(-1) sum_{i=1}^n [ G_n,P,u(O_i) - E_P G_n,P,u(O) ]$ \quad (\text{inputs: } n, G_n,P,u, u_0, P, O_1, ..., O_n)
\end{definition}

\begin{definition}[def:alpha-specific-ci-handle]\label{def:alpha-specific-ci-handle}
$C_n_alpha handle := For fixed alpha in (0,infinity) and delta in (0,1), define the scalar endpoint statistic T_n,alpha = rho_n(alpha)(hat_theta_L_n - theta_L). Let L_P,alpha denote the phase-specific scalar weak limit of T_n,alpha from thm:endpoint-limit-law: the regular-plus-boundary Gaussian scalar law for alpha > 1, the logarithmic transition Gaussian scalar law for alpha = 1, and the scalar projection of the infinitely divisible extreme-shell plus truncated-shell law for 0 < alpha < 1. Let hat L_n,alpha be a data-driven distribution estimator satisfying ass:fixed-alpha-calibrator for the finite-sample law of T_n,alpha. Define hat q_-(alpha) = inf{t : hat L_n,alpha((-infinity,t]) >= delta/2} and hat q_+(alpha) = inf{t : hat L_n,alpha((-infinity,t]) >= 1 - delta/2}. The fixed-alpha confidence interval is the explicit sample map C_n_alpha(O_1,...,O_n) = [hat_theta_L_n - rho_n(alpha)^(-1) hat q_+(alpha), hat_theta_L_n - rho_n(alpha)^(-1) hat q_-(alpha)].$ \quad (\text{inputs: } C_n_alpha, alpha, hat_theta_L_n, phi_P, R_nP, H_nP_u0, Sigma_P,u0, Lambda_P,alpha,u0, u_0)
\end{definition}

\section{Main results}

\begin{proposition}[prop:phase-normalization]\label{prop:phase-normalization}
Assume ass:iid and let P belong to P_n(alpha); among the class clauses, this calculation uses only ass:centered-shell-variance. Then there exist constants 0 < c_1 <= c_2 < infinity, depending only on the class constants, such that c_1 <= Var_P(S_nP) <= c_2 for all large n. Equivalently, the normalization rho_n(alpha) = sqrt(n / K_alpha(n^(-1))) is the unique order that keeps the boundary score nondegenerate, so rho_n(alpha) is of order n^(1/2) when alpha > 1, of order (n / log n)^(1/2) when alpha = 1, and of order n^(alpha / 2) when 0 < alpha < 1.
\end{proposition}
Proof. Because P belongs to P_n(alpha), def:boundary-class places the law in the main boundary class; for this proposition the only quantitative clause used from that class is ass:centered-shell-variance. Write f_n(o)=chi_P(o) B(o)^(-1/2) 1{n^(-1)<B(o)<=bar_b}. By def:boundary-score and def:shell-summand, S_nP=rho_n(alpha) n^(-1) sum_{i=1}^n {f_n(O_i)-P f_n}. By ass:iid, Var_P(S_nP)=rho_n(alpha)^2 n^(-1) Var_P(f_n(O)). By def:alpha-normalization this equals Var_P(f_n)/K_alpha(n^(-1)). Applying the centered-shell-variance clause from def:boundary-class at t=n^(-1) gives c_V_minus <= Var_P(S_nP) <= c_V_plus for all large n with n^(-1)<=bar_b. Thus take c_1=c_V_minus and c_2=c_V_plus. If a different deterministic normalization a_n were used in place of rho_n(alpha), the same computation would multiply the variance by (a_n/rho_n(alpha))^2; hence nondegeneracy bounded away from both zero and infinity is possible exactly when a_n is of the same order as rho_n(alpha). Finally def:tail-scale gives K_alpha(n^(-1))=1 for alpha>1, K_alpha(n^(-1))=1+log n for alpha=1, and K_alpha(n^(-1))=n^(1-alpha) for 0<alpha<1. Therefore rho_n(alpha) is respectively of order n^(1/2), (n/log n)^(1/2), and n^(alpha/2).
\par\smallskip\noindent \textit{Justification.} The proposal needs one explicit alpha-indexed normalization before any limit experiment or adaptation claim can be well-posed. \textit{Gap.} TanBlanchetSyrgkanis2026PRTEOT proves a continuous-instrument endpoint rate but does not identify the boundary-local normalization or the alpha = 1 transition. DaouiaFlorensSimar2010FrontierEVT supplies adjacent boundary-tail variance bookkeeping in a frontier problem, not for a CCOT PRTE endpoint with centered shell nondegeneracy. \textit{Consumer.} This proposition fixes the scaling used in thm:local-experiment, thm:endpoint-limit-law, and the fixed-alpha CI handle.

\begin{theorem}[thm:local-experiment]\label{thm:local-experiment}
Fix any u_0 in (0,1] and assume ass:iid. Let P belong to P_n(alpha). Assume additionally ass:shell-increment-moment; for alpha = 1 assume ass:boundary-annulus-oscillation for the nonextension-to-zero clause; and for 0 < alpha < 1 assume ass:stable-shell-drift and ass:stable-shell-nondegenerate-jump. Then the truncated boundary process H_nP_u0 is tight in D([u_0,1]) and converges weakly to a nondegenerate local limit experiment H_P,alpha,u0. If alpha > 1, H_P,alpha,u0 is a mean-zero Gaussian process with covariance kernel Sigma_P,u0. If alpha = 1, H_P,alpha,u0 is a logarithmic-transition Gaussian process with the same covariance kernel on every fixed truncated domain [u_0,1], but no extension to u_0 = 0 with the same normalization. If 0 < alpha < 1, the finite-dimensional truncated-shell marginals of the jointly infinitely divisible extreme-shell/truncated-shell vectors indexed by Lambda_P,alpha,u0 determine a unique continuous-path local limit experiment H_P,alpha,u0 on [u_0,1].
\end{theorem}
Proof. Because P belongs to P_n(alpha), def:boundary-class fixes the intended atom-free exact-tail active-face regime and also supplies the branch shell premise: ass:shell-row-lindeberg when alpha>=1 and ass:stable-shell-domain when 0<alpha<1. Fix u_0 in (0,1]. For finite u_1,...,u_m in [u_0,1], H_nP_u0(u_j)=sum_{i=1}^n X_{n,i,j} with X_{n,i,j}=n^(-1){G_n,P,u_j(O_i)-P G_n,P,u_j}. In the branch alpha>=1, lem:gaussian-shell-fidi gives convergence of every such vector to a centered Gaussian vector with covariance Sigma_P,u0(u_j,u_k). In the branch 0<alpha<1, projecting the joint convergence in lem:stable-shell-fidi onto the truncated-shell coordinates gives convergence of every finite vector (H_nP_u0(u_1),...,H_nP_u0(u_m)) to the corresponding truncated-shell marginal of the infinitely divisible law indexed by Lambda_P,alpha,u0. Lem:truncated-process-tightness supplies both tightness of H_nP_u0 in D([u_0,1]) and the continuity premise required by lem:process-convergence-from-fidi-and-tightness, namely that every subsequential weak limit has a continuous modification. Therefore lem:process-convergence-from-fidi-and-tightness upgrades the finite-dimensional limits in each branch to weak convergence in D([u_0,1]) toward a unique limit law H_P,alpha,u0. When alpha>1 this law is the mean-zero Gaussian process indexed by Sigma_P,u0. When alpha=1 the same fixed-domain Gaussian limit holds, but lem:alpha-one-no-zero-extension, using ass:boundary-annulus-oscillation, rules out an extension to u_0=0 with the same normalization and a right-continuous/tight origin. When 0<alpha<1, the limit process is the unique continuous-path law whose finite-dimensional marginals are the truncated-shell coordinates of the infinitely divisible extreme-shell/truncated-shell vectors coming from Lambda_P,alpha,u0. Nondegeneracy follows from lem:local-limit-nondegeneracy: for alpha>=1 the variance at u=1 is bounded away from zero by the centered-shell-variance clause inside def:boundary-class, while for 0<alpha<1 the jump projection is nonzero by ass:stable-shell-nondegenerate-jump.
\par\smallskip\noindent \textit{Justification.} The endpoint law is driven by a shell-local empirical process, so the proposal needs the underlying local experiment stated on a domain where the covariance or jump law is explicitly identified. \textit{Gap.} FangSantos2019DirectionalInference gives a generic directional-inference template and TanBlanchetSyrgkanis2026PRTEOT gives rate-only CCOT PRTE bounds, but neither paper derives the marked boundary-shell experiment for the active-face PRTE endpoint. The present statement adds the exact shell covariance or Levy object needed to make the phase-specific process claim well-posed. \textit{Consumer.} This theorem is the process-level input for endpoint calibration and for understanding why the alpha<=1 branch cannot be globalized to u=0 at the same normalization.

\begin{proposition}[prop:boundary-nonuniformity]\label{prop:boundary-nonuniformity}
Assume ass:iid and let P belong to P_n(alpha) with alpha <= 1. Assume additionally ass:shell-increment-moment; for alpha = 1 assume ass:boundary-annulus-oscillation; and for 0 < alpha < 1 assume ass:stable-shell-drift and ass:stable-shell-nondegenerate-jump. Then for every fixed u_0>0, H_nP_u0 is tight on [u_0,1] as in thm:local-experiment. The nonuniform phenomenon concerns only the limit u_0 downarrow 0. If alpha = 1, then for every sufficiently small delta_0>0 there exist sequences u_n,v_n in (0,delta_0] and constants eta,epsilon>0 such that liminf_n P_P(|H_nP_delta0(u_n)-H_nP_delta0(v_n)|>eta)>=epsilon; equivalently, the fixed-truncation family admits no tight right-continuous extension to u=0 with the same normalization. If 0<alpha<1, the explicit extreme-shell coordinate R_nP from 0 < B <= n^(-1) is first order in the joint extreme-shell/truncated-shell decomposition and is not o_P(1). No broader claim is made here about a generic scalar-endpoint omitted-shell remainder or a full-domain ell_infinity impossibility without extra u-indexed input.
\end{proposition}
Proof. The fixed-u_0 assertion is exactly the tightness part of thm:local-experiment. For alpha=1, lem:alpha-one-no-zero-extension is already formulated in terms of the typed truncated process H_nP_delta0 and gives, for every sufficiently small delta_0, sequences u_n,v_n in (0,delta_0] and constants eta,epsilon>0 such that liminf_n P_P(|H_nP_delta0(u_n)-H_nP_delta0(v_n)|>eta)>=epsilon. That is the displayed oscillation claim, and the same lemma gives the equivalent no-tight-zero-extension formulation. For 0<alpha<1, lem:extreme-shell-first-order applies the joint infinitely divisible convergence and ass:stable-shell-nondegenerate-jump to the first, extreme-shell coordinate. It yields convergence of R_nP to a nondegenerate infinitely divisible random variable, so R_nP is not o_P(1). This is precisely the claimed nonuniform phenomenon: the fixed truncation H_nP_u0 is tight for every u_0>0, but the omitted shell 0<B<=n^(-1) remains first order in the joint extreme-shell/truncated-shell decomposition. The statement makes no additional scalar-endpoint or full-domain ell_infinity claim beyond these two consequences.
\par\smallskip\noindent \textit{Justification.} The reviewer-tested fault line is exactly the behavior near u = 0, so the proposal should state the nonuniformity rather than silently overclaim a full-domain process. \textit{Gap.} DaouiaFlorensSimar2010FrontierEVT and KimPollard1990CubeRoot both show that boundary limits can fail to globalize naively, but neither identifies the alpha = 1 pointwise-versus-uniform split for a CCOT PRTE endpoint process. \textit{Consumer.} This proposition prevents invalid full-domain bootstrap or process-level CI claims and explains why oeq:boundary-honest-ci must separate fixed truncation from the u_0 downarrow 0 remainder.

\begin{theorem}[thm:endpoint-limit-law]\label{thm:endpoint-limit-law}
Assume ass:iid and let P belong to P_n(alpha); among the class clauses, the proof below uses only ass:boundary-atom-free quantitatively, while the remaining P_n(alpha) clauses keep the endpoint law scoped to the same active-face boundary regime as thm:local-experiment and oeq:boundary-honest-ci. Assume additionally ass:regular-influence-centering, ass:ccot-gamma-hadamard, and ass:ccot-donsker; for alpha >= 1 assume ass:extreme-shell-first-moment; for alpha > 1 assume ass:regular-boundary-joint-limit; and for 0 < alpha < 1 assume ass:shell-increment-moment, ass:stable-shell-drift, ass:stable-shell-nondegenerate-jump, and ass:irregular-endpoint-projection-nondegeneracy. Then the empirical CCOT plug-in endpoint admits the decomposition rho_n(alpha) { hat_theta_L_n - theta_L } = rho_n(alpha) n^(-1) sum_{i=1}^n phi_P(O_i) + S_nP + R_nP + o_P(1). Under the atom-free boundary clause inside P_n(alpha) there is no additional atom-at-B = 0 term beyond the explicit extreme-shell residual R_nP. Consequently, if alpha > 1 then R_nP = o_P(1), rho_n(alpha) = n^(1/2), and the limit law is Gaussian, combining the regular interior influence term with the Gaussian boundary factor indexed by Sigma_P,u0 and the regular-boundary cross-covariance. If alpha = 1 then R_nP = o_P(1), the regular interior term is asymptotically negligible, and the endpoint has a distinct logarithmic transition law. If 0 < alpha < 1, then for every fixed u_0 in (0,1] the pair (R_nP,H_nP_u0) converges jointly in R x D([u_0,1]) to the unique law whose finite-dimensional coordinates are the infinitely divisible vectors indexed by Lambda_P,alpha,u0, and the scalar endpoint projection R_nP+H_nP_u0(1) is boundary-irregular, nondegenerate, and non-Gaussian.
\end{theorem}
Proof. Because P belongs to P_n(alpha), def:boundary-class keeps the theorem inside the note's main active-face boundary regime; quantitatively, only its atom-free clause is used below. First apply lem:ccot-vonmises-expansion-from-delta. Since theta_L=Gamma(P) and hat_theta_L_n=Gamma(hat_P_n), it gives rho_n(alpha){hat_theta_L_n-theta_L}=rho_n(alpha)(hat_P_n-P)phi_P+S_nP+R_nP+o_P(1). Ass:regular-influence-centering gives P phi_P=0, so rho_n(alpha)(hat_P_n-P)phi_P=rho_n(alpha)n^(-1)sum_{i=1}^n phi_P(O_i). This proves the displayed decomposition. The same lemma already uses the atom-free boundary clause inside def:boundary-class to rule out any additional B=0 atom contribution beyond the explicit extreme-shell residual R_nP. If alpha>1, def:tail-scale gives rho_n(alpha)=n^(1/2). Lem:extreme-shell-negligible gives R_nP=o_P(1), and lem:regular-boundary-gaussian-sum applies ass:regular-influence-centering together with ass:regular-boundary-joint-limit to the joint row array of n^(-1/2)phi_P and the boundary score. Hence the scalar endpoint limit is Gaussian with the stated regular-boundary cross-covariance. If alpha=1, lem:extreme-shell-negligible again gives R_nP=o_P(1), while lem:regular-term-negligible-under-boundary-scaling gives rho_n(alpha)n^(-1)sum_i phi_P(O_i)=o_P(1). The remaining shell term is S_nP=H_nP_1(1), and lem:gaussian-shell-fidi at the single coordinate u_0=u_1=1 yields the logarithmic-transition Gaussian law under the normalization fixed in prop:phase-normalization. If 0<alpha<1, lem:regular-term-negligible-under-boundary-scaling removes the regular term. The m=0 case of lem:stable-shell-fidi gives tightness and convergence of the scalar extreme-shell coordinate R_nP, while lem:truncated-process-tightness gives tightness of H_nP_u0 and the continuity of its subsequential limits from ass:shell-increment-moment. Applying lem:joint-scalar-process-convergence-from-fidi-and-tightness to the joint finite-dimensional laws supplied by lem:stable-shell-fidi yields joint weak convergence of (R_nP,H_nP_u0) in R x D([u_0,1]) to the unique law indexed by Lambda_P,alpha,u0. Finally, lem:irregular-endpoint-projection-limit discharges the scalar endpoint conclusion by identifying the projection R_nP+H_nP_u0(1) as nondegenerate and non-Gaussian. Relative to the in-repository note stat_ate_overlap_decay, which develops overlap-tail phase bookkeeping for ATE risk rather than a CCOT endpoint law, the present contribution is the law-level phase transition for the sharp PRTE OT endpoint itself, conditional on the high-level active-face linearization together with the empirical-path admissibility input encoded in ass:ccot-gamma-hadamard and ass:ccot-donsker.
\par\smallskip\noindent \textit{Justification.} This is the main inference kernel once a local active-face expansion of the CCOT endpoint has been supplied: it isolates the sharp endpoint law and makes the alpha = 1 phase transition explicit rather than burying it inside a rate statement. \textit{Gap.} TanBlanchetSyrgkanis2026PRTEOT stops at convergence rates for the continuous-instrument endpoint, while FranguridiLiu2025RegularizedOTPartialID regularizes the OT value instead of handling the unregularized boundary endpoint. The closest in-repository comparator, stat_ate_overlap_decay, organizes overlap-tail phase behavior for ATE risk rather than for a sharp CCOT PRTE endpoint law. The present statement isolates the unregularized CCOT endpoint's regular-versus-irregular limit split, now with the omitted extreme shell made explicit rather than absorbed by fiat, but it does so conditionally on the high-level active-face linearization rather than deriving that singular derivative from CCOT primitives inside this core. \textit{Consumer.} It upgrades the anchor paper from rate-only asymptotics to a referee-checkable endpoint law, which is the missing prerequisite for any fixed-alpha CI theory.

\begin{proposition}[prop:strict-overlap-reduction]\label{prop:strict-overlap-reduction}
Assume ass:iid, ass:regular-influence-centering, and ass:strict-overlap-empirical-smoothness, and let P belong to P_n^so(p_so); quantitatively, only the ass:strict-overlap-boundary-regularity clause of the comparator class is used below, the remaining class clauses keep the strict-overlap geometry aligned with the main endpoint setup, and ass:regular-influence-centering is retained only to keep the comparator branch written with the same canonical phi_P convention as the main endpoint law. Then n^(1/2) { hat_theta_L_n - theta_L } = n^(-1/2) sum_{i=1}^n phi_P(O_i) + o_P(1). In particular, for every deterministic sequence t_n downarrow 0 with t_n <= bar_b, the centered root-n contribution of chi_P(O) B(O)^(-1/2) 1{ 0 < B(O) <= t_n } is o_P(1). No separate claim is made that a nonshrinking square-integrable boundary term vanishes on its own; any such contribution is absorbed into the high-level first-order expansion supplied by ass:strict-overlap-empirical-smoothness. Consequently the endpoint statistic has the regular centered Gaussian limit N(0,Var_P(phi_P(O))), allowing the variance-zero degenerate case.
\end{proposition}
Proof. Work under the assumptions of the proposition and fix P in the strict-overlap comparator class. By def:strict-overlap-class, this keeps the same strict-overlap geometry and endpoint labeling as the main branch; quantitatively, only its boundary-regularity clause is used below. Ass:regular-influence-centering is retained only to keep the canonical centered representative phi_P synchronized with the main endpoint-law branch; the actual comparator expansion comes from ass:strict-overlap-empirical-smoothness itself. The causal object is the lower sharp CCOT endpoint theta_L=Gamma(P), while the estimator is the plug-in map hat_theta_L_n=Gamma(hat_P_n). By lem:strict-overlap-plugin-vonmises, the strict-overlap empirical smoothness condition gives the first-order plug-in expansion
\[
 n^{1/2}\{\hat\theta_{L,n}-\theta_L\}
 = n^{-1/2}\sum_{i=1}^n \phi_P(O_i)+o_P(1).
\]
This proves the displayed root-n reduction.

Now let t_n downarrow 0 with t_n <= bar_b and set
\[
 g_{t_n}(O)=\chi_P(O)B(O)^{-1/2}{\bf 1}\{0<B(O)<=t_n\}.
\]
The shrinking-shell square-integrability clause encoded in def:strict-overlap-class and used by lem:shrinking-boundary-shell-negligible makes the boundary shell negligible, giving
\[
 n^{-1/2}\sum_{i=1}^n\{g_{t_n}(O_i)-E_P g_{t_n}(O)\}=o_P(1).
\]
This is exactly the asserted centered root-n negligibility of every deterministic shrinking boundary shell. No claim is made that a nonshrinking boundary term vanishes separately; such square-integrable first-order components are already absorbed into the regular influence representation supplied by the strict-overlap empirical smoothness condition.

Finally, lem:regular-rootn-gaussian-template yields
\[
 n^{-1/2}\sum_{i=1}^n\phi_P(O_i) \Rightarrow N\{0,\operatorname{Var}_P(\phi_P(O))\},
\]
with the variance-zero case interpreted as the degenerate Gaussian law. Slutsky's theorem applied to the first display gives
\[
 n^{1/2}\{\hat\theta_{L,n}-\theta_L\}\Rightarrow N\{0,\operatorname{Var}_P(\phi_P(O))\}.
\]
Thus this strict-overlap branch is only a sanity/comparator reduction: once the CCOT active-face boundary shell is square-integrable on shrinking neighborhoods and the endpoint is regular as a plug-in functional, the partial-identification endpoint collapses to the standard regular bound-estimator template.
\par\smallskip\noindent \textit{Justification.} A separate sanity branch is needed to show that the irregular alpha-tail machinery collapses to the familiar regular endpoint problem under genuine strict overlap. \textit{Gap.} ImbensManski2004PartialCI and Stoye2009MorePartialCI are the canonical regular-endpoint CI templates. This proposition is only a reduction result, but it is the exact comparator statement showing that the present kernel is not a relabeling of that regular case. \textit{Consumer.} It anchors the proposal to the standard partially identified CI benchmark and explains what disappears once the lower propensity is bounded away from the external boundary.

\begin{openendedquestion}[oeq:boundary-honest-ci]\label{oeq:boundary-honest-ci}
Using the scalar interval in def:alpha-specific-ci-handle and assuming the fixed-alpha law estimator satisfies ass:fixed-alpha-calibrator, one can construct intervals C_n_alpha such that liminf_{n -> infinity} inf_{P in P_n(alpha)} P_P( theta_L in C_n_alpha ) >= 1 - delta and sup_{P in P_n(alpha)} E_P[ |C_n_alpha| ] = O( rho_n(alpha)^(-1) ) for a fixed alpha, while remaining valid at the alpha = 1 transition.
\end{openendedquestion}
Proof. This is a conditional fixed-alpha result; the open object is the construction of a calibrator satisfying ass:fixed-alpha-calibrator from primitive estimators. Conditional on that assumption, def:alpha-specific-ci-handle defines T_n,alpha=rho_n(alpha)(hat_theta_L_n-theta_L), generalized quantiles hat q_-(alpha), hat q_+(alpha), and C_n_alpha=[hat_theta_L_n-rho_n(alpha)^(-1)hat q_+(alpha),hat_theta_L_n-rho_n(alpha)^(-1)hat q_-(alpha)]. The event theta_L in C_n_alpha is exactly {hat q_-(alpha)<=T_n,alpha<=hat q_+(alpha)}. Lem:fixed-alpha-ci-coverage uses ass:fixed-alpha-calibrator to show that this event has asymptotic probability at least 1-delta uniformly over P in P_n(alpha), including the alpha=1 transition because the calibrated law is the scalar logarithmic transition law from thm:endpoint-limit-law. Lem:fixed-alpha-ci-length gives |C_n_alpha|=rho_n(alpha)^(-1){hat q_+(alpha)-hat q_-(alpha)} and a uniform expected length O(rho_n(alpha)^(-1)). Combining the two lemmas proves the claimed conditional construction.
\par\smallskip\noindent \textit{Justification.} The endpoint law is only operational if the positive fixed-alpha CI problem is exposed honestly as the remaining construction task. \textit{Gap.} FangSantos2019DirectionalInference and FranguridiLiu2025RegularizedOTPartialID suggest derivative-based calibration routes for nonsmooth OT functionals, but neither paper constructs a fixed-alpha honest CI for the unregularized CCOT PRTE endpoint with the alpha = 1 transition handled explicitly. \textit{Consumer.} A positive answer would turn the limit-law kernel into an operational inference procedure instead of a pure lower-bound and asymptotic-calibration result.

\begin{lemma}[lem:lindeberg-feller-vector-array]\label{lem:lindeberg-feller-vector-array}
For a triangular array of independent mean-zero R^m-valued row variables whose row-sum covariance matrices converge to Sigma and whose array satisfies the multivariate Lindeberg condition, the row sums converge in distribution to N(0,Sigma).
\end{lemma}
Let X_{n,1},...,X_{n,k_n} be independent mean-zero random vectors in R^m. If sum_i Cov(X_{n,i}) -> Sigma and, for every epsilon>0, sum_i E[||X_{n,i}||^2 1{||X_{n,i}||>epsilon}] -> 0, then sum_i X_{n,i} converges weakly to the centered normal law N_m(0,Sigma).

\begin{lemma}[lem:kolmogorov-tightness-criterion]\label{lem:kolmogorov-tightness-criterion}
A sequence of processes on a compact interval with uniformly bounded q-th moment at one point and increment bound E|X_n(t)-X_n(s)|^q <= C|t-s|^beta for q>0 and beta>1 is tight in D on that interval, and every subsequential limit has a continuous modification.
\end{lemma}
If for some alpha>0, beta>1 and nondecreasing continuous F one has E|X_n(t)-X_n(s)|^alpha <= |F(t)-F(s)|^beta uniformly in n, together with tightness of X_n(t_0), then the laws of X_n are tight in D, and the modulus of continuity is tight.

\begin{lemma}[lem:triangular-array-id-convergence]\label{lem:triangular-array-id-convergence}
For a triangular array of independent infinitesimal row variables in R^m, vague convergence of n times the row law on sets bounded away from zero, convergence of compensated drift terms for a common truncation function, and negligible Gaussian part imply convergence of the row sum to the infinitely divisible law with the stated Levy measure and drift.
\end{lemma}
For an infinitesimal triangular array, convergence of the Levy measures on continuity sets bounded away from zero, convergence of the truncated first moments, and convergence of the small-jump covariance term characterize weak convergence of the row sums to the infinitely divisible distribution with that Levy triplet.

\begin{lemma}[lem:ccot-vonmises-via-delta-method]\label{lem:ccot-vonmises-via-delta-method}
Under ass:ccot-gamma-hadamard and ass:ccot-donsker, Gamma(hat_P_n)-Gamma(P)=(hat_P_n-P)phi_P+(hat_P_n-P)f_B+o_P(rho_n(alpha)^(-1)), where f_B(O)=chi_P(O)B(O)^(-1/2)1{0<B(O)<=bar_b}.
\end{lemma}
Let U_P, ||.||_F, f_B, and eta_P be the objects supplied by ass:ccot-gamma-hadamard. On the event {hat_P_n in U_P}, that assumption applied at Q=hat_P_n gives Gamma(hat_P_n)-Gamma(P)=(hat_P_n-P)phi_P+(hat_P_n-P)f_B+Rem_P(hat_P_n) with |Rem_P(hat_P_n)| <= ||hat_P_n-P||_F eta_P(||hat_P_n-P||_F). Ass:ccot-donsker gives P(hat_P_n in U_P)->1 and ||hat_P_n-P||_F=O_P(rho_n(alpha)^(-1)). Hence eta_P(||hat_P_n-P||_F)=o_P(1), so the remainder is o_P(rho_n(alpha)^(-1)). This yields the claimed empirical-path linearization.

\begin{lemma}[lem:gaussian-shell-fidi]\label{lem:gaussian-shell-fidi}
Under ass:iid and ass:shell-row-lindeberg, for alpha>=1 and fixed u_0 in (0,1], every finite vector (H_nP_u0(u_1),...,H_nP_u0(u_m)) with u_j in [u_0,1] converges to a centered Gaussian vector with covariance Sigma_P,u0(u_j,u_k).
\end{lemma}
For fixed u_1,...,u_m, def:truncated-boundary-process writes the vector as sum_{i=1}^n X_{n,i}, where X_{n,i,j}=n^(-1){G_n,P,u_j(O_i)-P G_n,P,u_j}. The variables are independent and mean zero by ass:iid. Ass:shell-row-lindeberg states exactly that this row array satisfies multivariate Lindeberg and that sum_i Cov(X_{n,i,j},X_{n,i,k})=n^(-1)Cov_P(G_n,P,u_j,G_n,P,u_k) converges to Sigma_P,u0(u_j,u_k). Lem:lindeberg-feller-vector-array gives the claimed Gaussian finite-dimensional convergence.

\begin{lemma}[lem:stable-shell-fidi]\label{lem:stable-shell-fidi}
For 0<alpha<1, under ass:iid, ass:stable-shell-domain, and ass:stable-shell-drift, the joint row sums (R_nP,H_nP_u0(u_1),...,H_nP_u0(u_m)) converge to the infinitely divisible law with Levy measure Lambda_P,alpha,u0 and the corresponding compensated drift.
\end{lemma}
By def:extreme-shell-residual and def:truncated-boundary-process, the joint vector is the centered row sum of n^(-1)(J_n,P,G_n,P,u_1,...,G_n,P,u_m). Ass:stable-shell-domain gives vague convergence of n times the row law on sets bounded away from zero to Lambda_P,alpha,u0 and negligible Gaussian part. Ass:stable-shell-drift gives convergence of the compensated drift terms for a common truncation function. With ass:iid, the row variables are independent and identically distributed within each row, so lem:triangular-array-id-convergence applies and yields the stated infinitely divisible finite-dimensional limit.

\begin{lemma}[lem:truncated-process-tightness]\label{lem:truncated-process-tightness}
Under ass:shell-increment-moment, H_nP_u0 is tight in D([u_0,1]) for every fixed u_0 in (0,1], and every subsequential weak limit has a continuous modification.
\end{lemma}
Ass:shell-increment-moment gives q>2, beta>1, and C<infinity such that sup_n E|H_nP_u0(v)-H_nP_u0(u)|^q <= C|v-u|^beta for all u,v in [u_0,1], and sup_n E|H_nP_u0(u_0)|^q <= C. The one-point laws are therefore tight by Markov's inequality, and the increment display is the hypothesis of lem:kolmogorov-tightness-criterion. Hence the laws of H_nP_u0 are tight in D([u_0,1]), and the same criterion yields that every subsequential weak limit has a continuous modification.

\begin{lemma}[lem:process-convergence-from-fidi-and-tightness]\label{lem:process-convergence-from-fidi-and-tightness}
Let X_n be D([u_0,1])-valued. If X_n is tight in D([u_0,1]), every subsequential limit has a continuous modification, and for every finite u_1,...,u_m in [u_0,1] the vector (X_n(u_1),...,X_n(u_m)) converges in distribution to a law nu_{u_1,...,u_m}, then there exists a unique probability law on C([u_0,1]) with those finite-dimensional marginals, viewed inside D([u_0,1]), and X_n converges weakly to that law.
\end{lemma}
Tightness gives subsequential weak limits in D([u_0,1]). Let X be one such limit. Because X has a continuous modification, the evaluation map x -> (x(u_1),...,x(u_m)) is continuous X-almost surely for every finite collection u_1,...,u_m, so the continuous-mapping theorem transfers the assumed finite-dimensional limits to X. Thus every subsequential limit has the same finite-dimensional distributions nu_{u_1,...,u_m}. A probability law on C([u_0,1]) is determined by its finite-dimensional distributions on a dense subset, so all subsequential limits coincide after passing to their continuous versions. Therefore the full sequence converges weakly in D([u_0,1]) to that unique continuous-path law.

\begin{lemma}[lem:joint-scalar-process-convergence-from-fidi-and-tightness]\label{lem:joint-scalar-process-convergence-from-fidi-and-tightness}
Let R_n be R-valued and X_n be D([u_0,1])-valued. If R_n is tight in R, X_n is tight in D([u_0,1]), every subsequential limit of X_n has a continuous modification, and for every finite u_1,...,u_m in [u_0,1] the vector (R_n,X_n(u_1),...,X_n(u_m)) converges in distribution to a law nu_{u_1,...,u_m}, then (R_n,X_n) converges weakly in R x D([u_0,1]) to the unique joint law with those finite-dimensional coordinates and continuous-path process coordinate.
\end{lemma}
Tightness of R_n and X_n implies tightness of the pair (R_n,X_n) in R x D([u_0,1]). Let (R,X) be any subsequential limit. Because X has a continuous modification, the map (r,x) -> (r,x(u_1),...,x(u_m)) is continuous at (R,X)-almost every point, so the continuous-mapping theorem identifies the finite-dimensional laws of (R,X) with the limiting family nu_{u_1,...,u_m}. A probability law on R x C([u_0,1]) is determined by the joint finite-dimensional distributions of the scalar coordinate and the process evaluated on a dense subset, so every subsequential limit yields the same joint law. Hence the whole pair sequence converges weakly in R x D([u_0,1]) to that unique law.

\begin{lemma}[lem:local-limit-nondegeneracy]\label{lem:local-limit-nondegeneracy}
The local limit experiment in thm:local-experiment is nondegenerate: for alpha>=1 the Gaussian limit has a coordinate with positive variance, and for 0<alpha<1 the jump limit has a nonzero extreme-shell projection.
\end{lemma}
For alpha>=1, the coordinate u=1 equals the boundary score S_nP by def:boundary-score and def:truncated-boundary-process. The variance bounds in ass:centered-shell-variance, applied at t=n^(-1), imply liminf_n Var_P(S_nP)>0 after the normalization used in S_nP. Therefore the Gaussian finite-dimensional limit from lem:gaussian-shell-fidi has a nonzero variance coordinate. For 0<alpha<1, lem:stable-shell-fidi identifies the joint infinitely divisible law, and ass:stable-shell-nondegenerate-jump states that the projection of Lambda_P,alpha,u0 onto the extreme-shell coordinate is nonzero. A nonzero Levy projection gives a nondegenerate coordinate, so the local experiment is nondegenerate.

\begin{lemma}[lem:alpha-one-no-zero-extension]\label{lem:alpha-one-no-zero-extension}
When alpha=1, ass:boundary-annulus-oscillation implies that the fixed-truncation Gaussian shell processes cannot be extended to u_0=0 with the same normalization as a tight right-continuous boundary process.
\end{lemma}
Ass:boundary-annulus-oscillation states that for every sufficiently small delta_0 there are u_n,v_n in (0,delta_0] and eta,epsilon>0 such that liminf_n P_P(|H_nP_delta0(u_n)-H_nP_delta0(v_n)|>eta)>=epsilon. If a tight right-continuous extension to u=0 with the same normalization existed, then the usual endpoint modulus condition for that extended family would force the oscillation over every sufficiently small truncation window to vanish, contradicting the displayed H_nP_delta0 oscillation. Hence no such zero-boundary extension exists.

\begin{lemma}[lem:extreme-shell-negligible]\label{lem:extreme-shell-negligible}
For alpha>=1, under ass:iid and ass:extreme-shell-first-moment, R_nP=o_P(1).
\end{lemma}
Let h_n(O)=chi_P(O)B(O)^(-1/2)1{0<B(O)<=n^(-1)}. By def:extreme-shell-residual, R_nP=n^(-1)sum_i rho_n(alpha){h_n(O_i)-P h_n}. Therefore E_P|R_nP| <= 2 rho_n(alpha) E_P|h_n(O)|. Ass:extreme-shell-first-moment gives E_P|h_n(O)| <= C n^{-(alpha-1/2)} for all large n. If alpha>1 then rho_n(alpha)=n^(1/2) by def:tail-scale and def:alpha-normalization, so the bound is O(n^(1-alpha))->0. If alpha=1 then rho_n(alpha)=(n/(1+log n))^(1/2), so the bound is O((1+log n)^(-1/2))->0. Markov's inequality gives R_nP=o_P(1).

\begin{lemma}[lem:regular-term-negligible-under-boundary-scaling]\label{lem:regular-term-negligible-under-boundary-scaling}
If alpha<=1 and ass:regular-influence-centering holds, then rho_n(alpha)n^(-1)sum_{i=1}^n phi_P(O_i)=o_P(1).
\end{lemma}
Ass:regular-influence-centering gives E_P phi_P=0 and E_P phi_P^2<infinity. By ass:iid, n^(-1/2)sum_i phi_P(O_i)=O_P(1), for example by Chebyshev's inequality. Thus rho_n(alpha)n^(-1)sum_i phi_P(O_i)=(rho_n(alpha)/sqrt(n)) O_P(1). If alpha=1 then rho_n(alpha)/sqrt(n)=(1+log n)^(-1/2)->0; if 0<alpha<1 then rho_n(alpha)/sqrt(n)=n^((alpha-1)/2)->0. Hence the regular term is o_P(1).

\begin{lemma}[lem:regular-boundary-gaussian-sum]\label{lem:regular-boundary-gaussian-sum}
For alpha>1, under ass:regular-influence-centering and ass:regular-boundary-joint-limit, the sum n^(-1/2)sum_i phi_P(O_i)+S_nP converges to the centered Gaussian law whose variance includes the regular variance, boundary variance, and regular-boundary cross-covariance.
\end{lemma}
For alpha>1, def:tail-scale and def:alpha-normalization give rho_n(alpha)=sqrt(n). Ass:regular-influence-centering supplies E_P phi_P(O)=0 and E_P phi_P(O)^2<infinity, so the first coordinate n^(-1/2)phi_P(O_i) is a legitimate mean-zero square-integrable row variable. The statistic n^(-1/2)sum_i phi_P(O_i)+S_nP is then a fixed linear combination of the joint triangular row sum formed by n^(-1/2)phi_P(O_i) and n^(-1)(G_n,P,1(O_i)-P G_n,P,1). Ass:regular-boundary-joint-limit states multivariate Lindeberg and convergence of the full covariance matrix, including cross-covariances. Lem:lindeberg-feller-vector-array gives joint Gaussian convergence, and applying the continuous linear map that sums the two coordinates gives the asserted Gaussian scalar limit.

\begin{lemma}[lem:ccot-vonmises-expansion-from-delta]\label{lem:ccot-vonmises-expansion-from-delta}
Under ass:ccot-gamma-hadamard, ass:ccot-donsker, and ass:boundary-atom-free, rho_n(alpha){Gamma(hat_P_n)-Gamma(P)}=rho_n(alpha)(hat_P_n-P)phi_P+S_nP+R_nP+o_P(1).
\end{lemma}
Let f_B=chi_P B^(-1/2)1{0<B<=bar_b}. Lem:ccot-vonmises-via-delta-method gives Gamma(hat_P_n)-Gamma(P)=(hat_P_n-P)phi_P+(hat_P_n-P)f_B+o_P(rho_n(alpha)^(-1)). Multiplying by rho_n(alpha) gives rho_n(alpha){Gamma(hat_P_n)-Gamma(P)}=rho_n(alpha)(hat_P_n-P)phi_P+rho_n(alpha)(hat_P_n-P)f_B+o_P(1). By ass:boundary-atom-free, f_B splits P-almost surely into chi_P B^(-1/2)1{0<B<=n^(-1)} plus chi_P B^(-1/2)1{n^(-1)<B<=bar_b}. The empirical fluctuation of the first piece multiplied by rho_n(alpha) is R_nP by def:extreme-shell-residual; the second is S_nP by def:boundary-score. This proves the display.

\begin{lemma}[lem:irregular-endpoint-projection-limit]\label{lem:irregular-endpoint-projection-limit}
For 0<alpha<1, the scalar projection R_nP+H_nP_u0(1) converges to a nondegenerate non-Gaussian infinitely divisible law under ass:irregular-endpoint-projection-nondegeneracy.
\end{lemma}
Taking u=1 in lem:stable-shell-fidi gives joint convergence of (R_nP,H_nP_u0(1)) to the corresponding two-coordinate infinitely divisible limit. The map (r,h)->r+h is continuous, so the scalar projection converges to the projected infinitely divisible law. Ass:irregular-endpoint-projection-nondegeneracy says that, after deterministic centering, this projected Levy measure is nonzero. A one-dimensional infinitely divisible law with nonzero Levy measure is non-Gaussian and nondegenerate. Therefore R_nP+H_nP_u0(1) has the asserted boundary-irregular scalar limit.

\begin{lemma}[lem:extreme-shell-first-order]\label{lem:extreme-shell-first-order}
For 0<alpha<1, under ass:stable-shell-nondegenerate-jump, the extreme-shell statistic R_nP is not o_P(1).
\end{lemma}
The first coordinate in lem:stable-shell-fidi is R_nP. Ass:stable-shell-nondegenerate-jump states that the Levy measure projected onto this coordinate is nonzero. Thus R_nP converges along the finite-dimensional limit to a nondegenerate infinitely divisible random variable. A sequence converging to a nondegenerate law cannot converge to zero in probability, so R_nP is not o_P(1).

\begin{lemma}[lem:fixed-alpha-ci-coverage]\label{lem:fixed-alpha-ci-coverage}
Conditional on ass:fixed-alpha-calibrator, the interval C_n_alpha in def:alpha-specific-ci-handle has asymptotic coverage at least 1-delta uniformly over P in P_n(alpha).
\end{lemma}
By def:alpha-specific-ci-handle, theta_L in C_n_alpha if and only if hat q_-(alpha)<=T_n,alpha<=hat q_+(alpha), where T_n,alpha=rho_n(alpha)(hat_theta_L_n-theta_L). Ass:fixed-alpha-calibrator states that the estimated distribution hat L_n,alpha uniformly approximates the finite-sample law of T_n,alpha at continuity points of the limiting law and that its delta/2 and 1-delta/2 quantiles are uniformly consistent. Therefore uniformly over P in P_n(alpha), P_P(T_n,alpha<hat q_-(alpha))<=delta/2+o(1) and P_P(T_n,alpha>hat q_+(alpha))<=delta/2+o(1). The union bound gives inf_P P_P(theta_L in C_n_alpha)>=1-delta-o(1). The phase-specific limiting laws whose continuity points are used here are those identified in thm:endpoint-limit-law, including the alpha=1 logarithmic transition law.

\begin{lemma}[lem:fixed-alpha-ci-length]\label{lem:fixed-alpha-ci-length}
Conditional on ass:fixed-alpha-calibrator, the interval C_n_alpha in def:alpha-specific-ci-handle satisfies sup_{P in P_n(alpha)} E_P|C_n_alpha|=O(rho_n(alpha)^(-1)).
\end{lemma}
The interval definition gives |C_n_alpha|=rho_n(alpha)^(-1){hat q_+(alpha)-hat q_-(alpha)}. Ass:fixed-alpha-calibrator includes the uniform bound E_P[hat q_+(alpha)-hat q_-(alpha)]=O(1). Multiplying by rho_n(alpha)^(-1) yields sup_{P in P_n(alpha)} E_P|C_n_alpha|=O(rho_n(alpha)^(-1)).

\begin{lemma}[lem:strict-overlap-neighborhood-entry]\label{lem:strict-overlap-neighborhood-entry}
Under ass:strict-overlap-empirical-smoothness and def:empirical-law, if N_P is the strict-overlap neighborhood of P on which the displayed local modulus bound holds, then P_P(hat_P_n in N_P) -> 1.
\end{lemma}
Proof. In ass:strict-overlap-empirical-smoothness, N_P is a local strict-overlap neighborhood of P for the seminorm ||Q-P||_{F_P}=sup_{f in F_P}|(Q-P)f|. Hence there is an epsilon_P>0 such that ||Q-P||_{F_P}<epsilon_P implies Q belongs to N_P. The same assumption gives ||hat_P_n-P||_{F_P}=O_P(n^{-1/2}), and therefore ||hat_P_n-P||_{F_P}=o_P(1). Consequently P_P(||hat_P_n-P||_{F_P}<epsilon_P)->1. On this event hat_P_n lies in N_P, so P_P(hat_P_n in N_P)->1.

\begin{lemma}[lem:strict-overlap-plugin-vonmises]\label{lem:strict-overlap-plugin-vonmises}
Under ass:iid, ass:strict-overlap-empirical-smoothness, def:strict-overlap-class, def:empirical-law, and def:plugin-endpoint, the plug-in endpoint satisfies n^(1/2){hat_theta_L_n-theta_L}=n^(-1/2) sum_{i=1}^n phi_P(O_i)+o_P(1).
\end{lemma}
Proof. Let N_P, F_P, ||.||_{F_P}, and eta_P be the objects supplied by ass:strict-overlap-empirical-smoothness for the law P in the strict-overlap class. By lem:strict-overlap-neighborhood-entry, P_P(hat_P_n in N_P)->1. On this event the local expansion in ass:strict-overlap-empirical-smoothness with Q=hat_P_n gives
\[
 \Gamma(\hat P_n)-\Gamma(P)=(\hat P_n-P)\phi_P+R_n,
 \qquad |R_n|\le ||\hat P_n-P||_{F_P}\eta_P(||\hat P_n-P||_{F_P}).
\]
The same assumption gives ||hat_P_n-P||_{F_P}=O_P(n^{-1/2}). Hence ||hat_P_n-P||_{F_P}=o_P(1), eta_P(||hat_P_n-P||_{F_P})=o_P(1), and
\[
 n^{1/2}|R_n|\le n^{1/2}||\hat P_n-P||_{F_P}\eta_P(||\hat P_n-P||_{F_P})=O_P(1)o_P(1)=o_P(1).
\]
By def:empirical-law,
\[
 (\hat P_n-P)\phi_P=n^{-1}\sum_{i=1}^n\phi_P(O_i)-E_P\phi_P(O).
\]
Assumption ass:strict-overlap-empirical-smoothness includes the centering condition E_P\phi_P(O)=0. Using def:plugin-endpoint, hat_theta_L_n=Gamma(hat_P_n), and theta_L=Gamma(P), multiplication by n^{1/2} gives the asserted expansion.

\begin{lemma}[lem:shrinking-boundary-shell-negligible]\label{lem:shrinking-boundary-shell-negligible}
Under ass:iid and ass:strict-overlap-boundary-regularity, for any deterministic sequence t_n downarrow 0 with t_n <= bar_b, if g_t(O)=chi_P(O)B(O)^(-1/2)1{0<B(O)<=t}, then n^(-1/2) sum_{i=1}^n {g_{t_n}(O_i)-E_P g_{t_n}(O)}=o_P(1).
\end{lemma}
Proof. By ass:strict-overlap-boundary-regularity with t=bar_b,
\[
 E_P[\chi_P(O)^2B(O)^{-1}{\bf 1}\{0<B(O)<=bar_b\}]<\infty.
\]
For t_n downarrow 0, the nonnegative variables chi_P(O)^2B(O)^{-1}1{0<B(O)<=t_n} converge pointwise to zero and are dominated by the integrable bar_b-shell envelope in the preceding display. Dominated convergence gives E_P[g_{t_n}(O)^2]->0. Under ass:iid, the variance of
\[
 n^{-1/2}\sum_{i=1}^n\{g_{t_n}(O_i)-E_Pg_{t_n}(O)\}
\]
equals Var_P(g_{t_n}(O)) and is bounded by E_P[g_{t_n}(O)^2]. This variance tends to zero, so Chebyshev's inequality implies convergence to zero in probability.

\begin{lemma}[lem:regular-rootn-gaussian-template]\label{lem:regular-rootn-gaussian-template}
Under ass:iid and ass:strict-overlap-empirical-smoothness, n^(-1/2) sum_{i=1}^n phi_P(O_i) converges in distribution to a centered Gaussian law with variance Var_P(phi_P(O)), allowing the variance-zero degenerate case.
\end{lemma}
Proof. Assumption ass:strict-overlap-empirical-smoothness includes E_P phi_P(O)=0 and E_P phi_P(O)^2<infinity. Under ass:iid, the variables phi_P(O_i) are iid with mean zero and finite variance sigma_P^2=Var_P(phi_P(O)). If sigma_P^2>0, the Lindeberg-Levy central limit theorem gives
\[
 n^{-1/2}\sum_{i=1}^n\phi_P(O_i)\Rightarrow N(0,sigma_P^2).
\]
If sigma_P^2=0, then phi_P(O)=0 P-almost surely after centering, so the normalized sum is identically zero and the limit is the degenerate Gaussian N(0,0).

\section{Honest scope}

The settled formal objects in this version are the alpha-indexed normalization rho_n(alpha), the marked boundary-shell regime with P(B = 0) = 0 and exact regular variation at the lower propensity edge, the truncated shell process H_nP_u0 on domains bounded away from u = 0, the explicit extreme-shell residual R_nP from 0 < B <= n^(-1), the typed nonuniformity statement showing why alpha <= 1 cannot support a tight right-continuous zero-extension at the same normalization, the endpoint decomposition and phase-specific limit law for the empirical CCOT plug-in endpoint, and the separate strict-overlap comparator class based on an external lower propensity floor. All headline claims remain conditional on the unique-active-face, separated-dual-gap subclass, on the high-level shell-limit conditions that identify Sigma_P,u0 or the joint Levy object Lambda_P,alpha,u0, and now explicitly on the provisional active-face local linearization / empirical admissibility pair encoded in ass:ccot-gamma-hadamard and ass:ccot-donsker. The open object is the positive fixed-alpha CI construction in oeq:boundary-honest-ci: the present core names the studentized local-experiment route but does not claim a completed honest interval. No cross-alpha honest-adaptation theorem is claimed in this version.

\begin{thebibliography}{99}

  \bibitem{TanBlanchetSyrgkanis2026PRTEOT} Tan, J., Blanchet, J., and Syrgkanis, V. (2026). Policy-Relevant Treatment Effects with Instruments via Optimal Transport. arXiv:2604.12263.

  \bibitem{SasakiUra2018PRTEInference} Sasaki, Y., and Ura, T. (2018). Estimation and Inference of Policy Relevant Treatment Effects. arXiv:1805.11503.

  \bibitem{DaouiaFlorensSimar2010FrontierEVT} Daouia, A., Florens, J.-P., and Simar, L. (2010). Frontier Estimation and Extreme Value Theory. arXiv:1011.5722.

  \bibitem{KimPollard1990CubeRoot} Kim, J., and Pollard, D. (1990). Cube Root Asymptotics. Annals of Statistics 18(1): 191-219.

  \bibitem{FangSantos2019DirectionalInference} Fang, Z., and Santos, A. (2019). Inference on Directionally Differentiable Functions. Econometric Theory 35(1): 1-38.

  \bibitem{FranguridiLiu2025RegularizedOTPartialID} Franguridi, G., and Liu, L. (2025). Regularized Optimal Transport Inference for Partially Identified Models. arXiv:2512.18084.

  \bibitem{ImbensManski2004PartialCI} Imbens, G. W., and Manski, C. F. (2004). Confidence Intervals for Partially Identified Parameters. Econometrica 72(6): 1845-1857.

  \bibitem{Stoye2009MorePartialCI} Stoye, J. (2009). More on Confidence Intervals for Partially Identified Parameters. Econometrica 77(4): 1299-1315.

  \bibitem{CaiLow2005AdaptiveCI} Cai, T. T., and Low, M. G. (2005). Adaptive Confidence Intervals for Nonparametric Function Estimation. Annals of Statistics 32(5): 2024-2047.

  \bibitem{RobinsVanderVaart2006AdaptiveSets} Robins, J., and van der Vaart, A. (2006). Adaptive Nonparametric Confidence Sets. Annals of Statistics 34(1): 229-253.

  \bibitem{ChernozhukovEtAl2018DML} Chernozhukov, V., Chetverikov, D., Demirer, M., Duflo, E., Hansen, C., Newey, W., and Robins, J. (2018). Double/Debiased Machine Learning for Treatment and Structural Parameters. Econometrics Journal 21(1): C1-C68.

  \bibitem{vanderVaart1998} van der Vaart, A. W. (1998). Asymptotic Statistics. Cambridge University Press (Ch. 20, functional delta method; Thm 20.8).

  \bibitem{Billingsley1999} Billingsley, P. (1999). Convergence of Probability Measures, 2nd ed. Wiley. (Tightness / Kolmogorov criterion.)

  \bibitem{Kallenberg2002} Kallenberg, O. (2002). Foundations of Modern Probability, 2nd ed. Springer. (Triangular-array / infinitely-divisible convergence.)

\end{thebibliography}

\end{document}